\chapter{Описание решения}
\label{cha:Solution}

В данной главе описываются реализованные алгоритмы укладки коробок на паллету.

\section{Структура HeightHandler}

Ключевым компонентом решения является структура \texttt{HeightHandler}, которая эффективно хранит информацию о текущих высотах на паллете.

\subsection{Принцип работы}

HeightHandler хранит множество непересекающихся прямоугольников с высотами. Каждый прямоугольник описывается координатами $(x, y, X, Y)$ и высотой $h$. При добавлении новой коробки:

\begin{enumerate}
    \item Существующие прямоугольники, пересекающиеся с новой коробкой, разрезаются на части
    \item Новый прямоугольник (с высотой = текущая высота + высота коробки) добавляется в множество
    \item Прямоугольники сортируются по убыванию высоты для быстрого поиска
\end{enumerate}

\subsection{Основные операции}

\begin{itemize}
    \item \texttt{get(x, y, X, Y)} --- возвращает максимальную высоту в прямоугольной области
    \item \texttt{add\_rect(rect)} --- добавляет новый прямоугольник, обновляя карту высот
    \item \texttt{get\_dots(header, box)} --- возвращает ``интересные'' точки для размещения коробки заданного размера
\end{itemize}

Функция \texttt{get\_dots} генерирует кандидатные позиции на основе углов существующих прямоугольников, что значительно сокращает пространство поиска.

\section{Жадный алгоритм (GreedySolver)}

GreedySolver реализует простой, но эффективный жадный подход.

\subsection{Алгоритм}

\begin{enumerate}
    \item Коробки предварительно сортируются по убыванию площади основания
    \item Для каждой коробки в отсортированном порядке:
    \begin{enumerate}
        \item Генерируются все возможные ориентации коробки (до 6 вариантов)
        \item Для каждой ориентации получаем кандидатные позиции из HeightHandler
        \item Выбирается позиция и ориентация с минимальной итоговой высотой $h + h_{box}$
        \item Коробка размещается, HeightHandler обновляется
    \end{enumerate}
\end{enumerate}

\subsection{Сложность}

Временная сложность: $O(n \cdot k \cdot m)$, где $n$ --- число коробок, $k$ --- число ориентаций (до 6), $m$ --- число кандидатных позиций (растёт с числом уложенных коробок).

\subsection{Результаты}

На тестовом наборе из 436 задач GreedySolver показывает:
\begin{itemize}
    \item Среднее заполнение: 73.56\%
    \item Минимум: 59.01\%
    \item Максимум: 85.21\%
    \item Время: 308 мс (на все 436 тестов)
\end{itemize}

\section{LNS с имитацией отжига (LNSSolver)}

LNSSolver использует метаэвристический подход Large Neighborhood Search (LNS) в сочетании с имитацией отжига (Simulated Annealing).

\subsection{Представление решения}

Решение кодируется как массив структур \texttt{BoxMeta}:
\begin{itemize}
    \item \texttt{box\_id} --- индекс коробки
    \item \texttt{rotation} --- выбранная ориентация (-1 = рассмотреть все)
    \item \texttt{position} --- предпочтительный угол паллеты (0-3)
\end{itemize}

Порядок элементов в массиве определяет порядок укладки коробок.

\subsection{Функция симуляции}

Функция \texttt{simulate} принимает массив BoxMeta и симулирует укладку:
\begin{enumerate}
    \item Для каждой коробки в заданном порядке:
    \begin{enumerate}
        \item Если rotation задан --- используем его, иначе перебираем все ориентации
        \item Выбираем позицию с минимальной высотой
        \item При равной высоте предпочитаем позицию ближе к указанному углу (position)
    \end{enumerate}
    \item Возвращается итоговая высота паллеты
\end{enumerate}

\subsection{Операторы мутации}

Применяются следующие операторы с заданными вероятностями:
\begin{itemize}
    \item Изменение position или rotation случайной коробки (вес 20)
    \item Реверс подмассива (вес 10)
    \item Shuffle подмассива (вес 20)
    \item Swap двух коробок (вес 10)
    \item Swap соседних коробок (вес 20)
\end{itemize}

\subsection{Критерий принятия}

Новое решение принимается если:
$$score_{new} < score_{best}$$
или с вероятностью:
$$P = \exp\left(\frac{score_{best} - score_{new}}{T}\right)$$

где $T$ --- температура, убывающая по формуле $T_{new} = T \cdot 0.999$ после каждой итерации.

При отсутствии улучшений более 1 секунды температура сбрасывается к 1.

\subsection{Результаты}

\begin{table}[ht]
\centering
\caption{Результаты LNSSolver при разном времени работы}
\begin{tabular}{|l|c|c|c|c|}
\hline
\textbf{Время} & \textbf{Среднее} & \textbf{Мин} & \textbf{Макс} & \textbf{Общее время} \\
\hline
1 сек & 77.18\% & 72.86\% & 86.39\% & 14 сек \\
5 сек & 78.23\% & 73.11\% & 88.15\% & 70 сек \\
30 сек & 78.93\% & 75.36\% & 88.65\% & 7 мин \\
300 сек & 79.71\% & 76.03\% & 90.39\% & 70 мин \\
\hline
\end{tabular}
\end{table}

\section{Визуализатор}

Для анализа результатов разработан Python-скрипт \texttt{visualizer.py}, генерирующий интерактивную 3D-визуализацию.

\subsection{Возможности}

\begin{itemize}
    \item Интерактивное вращение и масштабирование 3D-модели
    \item Пошаговый просмотр укладки (клавиши $\leftarrow$/$\rightarrow$ или A/D)
    \item Раскраска коробок по SKU или другим атрибутам
    \item Настройка прозрачности и отображения рёбер
\end{itemize}

\subsection{Использование}

\begin{verbatim}
python3 src/scripts/visualizer.py answers/261.csv \
    --pallet-length 1200 --pallet-width 800 \
    --out boxes.html --color-by sku \
    --initial-count 100000000 --opacity 0.4
\end{verbatim}

\section{Архитектура программы}

Программа реализована на C++20 со следующей структурой:

\begin{itemize}
    \item \texttt{src/objects/} --- структуры данных (Box, TestData, Answer, Metrics)
    \item \texttt{src/solvers/} --- алгоритмы решения
    \begin{itemize}
        \item \texttt{height\_handler.*} --- HeightHandler
        \item \texttt{greedy/greedy\_solver.*} --- GreedySolver
        \item \texttt{lns/lns\_solver.*} --- LNSSolver
    \end{itemize}
    \item \texttt{src/utils/} --- вспомогательные утилиты (время, рандом, assert)
    \item \texttt{src/scripts/} --- Python-скрипты визуализации
\end{itemize}

Сборка осуществляется через CMake с поддержкой многопоточности:
\begin{verbatim}
cmake -DTHREADS_NUM=32 ..
make
\end{verbatim}
